\section{Przedmiot badań}

Przedmiotem badań jest obciążenie kognitywne programisty podczas wprowadzania zmian w istniejącym kodzie programu
napisanego w języku Java.
Obciążenie kognitywne w tym badaniu mierzone jest pośrednio poprzez czas wymagany na zrozumienie i modyfikacje kodu.

W szczególności badanie koncentruje się na tym, czy i w jakim stopniu wysoki poziom sprzężenia utrudnia modyfikację
kodu.

Celem badawczym jest określenie wpływu stopnia sprzężenia istniejących klas na czas wymagany do wprowadzenia zmian w
kodzie przez programistę.

Następny podrozdział formułuje hipotezę badawczą, którą eksperyment ma za zadanie zweryfikować.

\section{Hipoteza badawcza}

Hipoteza badawcza zakłada, że jeśli klasa jest powiązana z inną klasą o znacząco różnym stopniu sprzężenia, to czas
potrzebny na wprowadzenie zmian w obszarze tej relacji wykazuje statystycznie istotną różnicę.

Weryfikacja tej hipotezy polega na zmierzeniu czasu wykonywania tego samego zadania przez programistów na dwóch
wariantach tego samego programu.
Jeden z nich będzie posiadał zmodyfikowany kod źródłowy tak, żeby zwiększyć stopień sprzężenia pomiędzy
zaimplementowanymi wewnątrz klasami.
Mierząc czas wykonania zadania, można zbadać statystycznie istotną różnicę w średnim czasie potrzebnym do
wprowadzenia zmian.

\section{Kod źródłowy do badania}

Istotną częścią badania jest poprawna manipulacja zmienną niezależną --- współczynnikiem sprzężenia klas.
Oba warianty programu rozwiązują ten sam problem, to znaczy, że są testowane przy użyciu tego samego zestawu
testów jednostkowych.

Ze względu na ograniczenia czasowe eksperyment nie mógł trwać dłużej niż godzinę dla każdego uczestnika.
Wyzwaniem postawionym przed autorem pracy było utworzenie kodu źródłowego, który jest:

\begin{itemize}
    \item niewielki rozmiarem, zawiera nie więcej niż 10 klas, a żaden z plików klas nie jest dłuższy niż 100 linii kodu
    \item zrozumiale napisany, nie wymaga dodatkowego wyjaśniania zaimplementowanych algorytmów czy wprowadzonych
    abstrakcyjnych pojęć
    \item na tyle skomplikowany, by zmodyfikować relację między klasami i nie naruszyć ogólnej struktury programu i
    rozwiązywanego problemu
    \item kompletnie weryfikowalny za pomocą przygotowanych testów jednostkowych
\end{itemize}

Kod źródłowy programu implementuje mechanizm magistrali wiadomości (ang. \textit{message bus}).
Jest to kategoria problemu powszechnie znana wśród programistów Java, nie wymaga znajomości dziedziny spoza
informatyki, a większość badanych od razu porównywała rozwiązywaną problematykę do programu \texttt{Kafka} — programu
realizującego magistralę danych w złożonym środowisku rozproszonym.

Tutaj powinienem przestawić diagram UML programu.

Zadaniem programu \texttt{JavaBus} jest przesyłanie wiadomości do uprzednio zarejestrowanych odbiorców.
Klasa \texttt{InMemoryBus} implementująca interfejs \texttt{Bus} jest głównym elementem programu.
Umożliwia ona zarejestrowanie klas implementujących interfejs \texttt{Consumer}, których metoda \texttt{consume}
zostanie wywołana w momencie wysłania obiektu implementującego klasę \texttt{Message} z odpowiadającym obiektem klasy
\texttt{Topic}.


\section{Opis eksperymentu}

Eksperyment obejmuje trzy kolejne etapy:

\begin{enumerate}
    \item wprowadzenia, w którym prowadzący badanie przedstawia kod źródłowy programu i problem, który rozwiązuje
    \item rozwiązania zadania pierwszego, w obszarze wspólnym dla obu grup badanych
    \item rozwiązania zadania drugiego, w obszarze o zmodyfikowanym współczynniku sprzężenia pomiędzy klasami
\end{enumerate}

Poniżej każdy etap jest szczegółowo opisany, wskazując jego cel i interakcje pomiędzy prowadzącym badanie a
badanym programistą.

\subsection{Wprowadzenie do badania}

Celem tego etapu jest wyjaśnienie programu i jego struktury badanemu programiście.
Badany nie powinien mieć problemu z określeniem rozwiązywanego problemu i potrafić wskazać elementy kodu, które wykonują
rejestrację konsumenta i przekazywanie wiadomości.
Jednocześnie wyjaśnienie powinno zawierać tę samą ilość informacji we wszystkich przeprowadzonych eksperymentach.

Badany programista otrzymuje kod źródłowy, z którym zapoznaje się w trakcie eksperymentu.
Badacz opisuje problem rozwiązywany przez program --- przekazywanie wiadomości zarejestrowanym konsumentom podkreślając,
że na potrzeby eksperymentu program zakłada funkcjonowanie wewnątrz jednego wątku procesu JVM, więc
programista nie musi rozważać współbieżności.

Na tym etapie programista może zadawać pytania co do przedstawionego problemu i zaprezentowanych klas.
Prowadzący badanie powinien odpowiedzieć wyłącznie na pytania, które odnoszą się do zaprezentowanego wyjaśnienia.

\subsection{Zadanie pierwsze}

Celem zadania pierwszego jest pomiar czasu zadania w obszarze kodu wspólnego dla obu wariantów.
Jednocześnie programista zapoznaje się z ogólną strukturą programu, by w zadaniu drugim skupić się na klasach o
zmodyfikowanym stopniu sprzężenia.

Prowadzący badanie instruuje uczestnika, aby odblokował pierwszy test, a następnie go uruchomił.
Po stwierdzeniu, że test kończy się błędem asercji, programista otrzymuje zadanie zmiany kodu.
Polega ona na wprowadzeniu obsługi klas niebezpośrednio implementujących \texttt{Message}.
Badacz informuje programistę, że może zmieniać wyłącznie implementacje programu, a interfejsy i testy musi pozostawić
nienaruszone.
Programista może wielokrotnie uruchamiać testy jednostkowe.
Badacz rejestruje czas rozpoczęcia zadania po wyjaśnieniu zadania, a czas zakończenia w chwili, gdy wszystkie testy
jednostkowe zakończą się sukcesem.
Jeżeli programista nie uruchomi testów jednostkowych po modyfikacji kodu, badacz poleca ich uruchomienie.

\subsection{Zadanie drugie}

Celem zadania drugiego jest właściwy pomiar czasu wprowadzenia zmiany w obszarze o różnym stopniu sprzężenia klas.
Programista powinien w tym czasie przeanalizować relacje pomiędzy klasami i wprowadzić oczekiwane zmiany.

Prowadzący badanie wyjaśnia, że zadanie drugie jest analogiczne do pierwszego.
Polega ono na zaimplementowaniu dodatkowej strategii informowania konsumentów o nowej wiadomości.
Klasa konsumenta może otrzymać wiadomość wyłącznie w sytuacji, gdy wszystkie uprzednio zarejestrowane obiekty klasy
\texttt{Consumer} zwróciły błąd po wywołaniu metody \texttt{consume}.
Programista odblokowuje dwa kolejne testy, uruchamia je i przystępuje do realizacji zadania.
Podobnie jak w zadaniu pierwszym, badacz rejestruje moment rozpoczęcia zadania i pomyślnego przetestowania
zmodyfikowanego kodu.